\documentclass[12pt,a4paper]{article}
\usepackage[utf8]{inputenc}
\usepackage[spanish]{babel}
\usepackage{amsmath, amssymb, geometry, enumitem, hyperref}

% Configuración de márgenes
\geometry{top=2.5cm, bottom=2.5cm, left=2.5cm, right=2.5cm}

% Configuración para los links
\hypersetup{
    colorlinks=true,
    linkcolor=blue,
    filecolor=magenta,      
    urlcolor=blue,
}

\title{\textbf{Informe Técnico T3: Aplicaciones de Probabilidad Conjunta}}
\author{Luis Rojas - 20222020242 \\ Asignatura: Probabilidad y Estadística}
\date{}

\begin{document}

\maketitle

\section*{Enlace de Validación en Python}
El desarrollo computacional y la validación de los resultados presentados en este informe se encuentran disponibles en el siguiente cuaderno de Google Colab: \\
\url{https://colab.research.google.com/drive/15V55QSXwe2XufjocfcXVe2EqaMGhKwzS?usp=sharing}

\section{Aplicación 1: Caso Discreto (Selección de Estudiantes)}
\textbf{Enunciado:} Se seleccionan 2 estudiantes de un grupo de 8 (3 de Sistemas, 2 de Electrónica y 3 de Industrial). Sea $X$ el número de estudiantes de Sistemas y $Y$ el de Electrónica. Se requiere calcular la probabilidad conjunta $P(X \le 1, Y \le 1)$.

\textbf{Fórmula de la Función de Probabilidad Conjunta:}
\[ f(x,y) = \frac{\binom{3}{x} \binom{2}{y} \binom{3}{2-x-y}}{\binom{8}{2}} \]

\textbf{Desarrollo del Denominador (Espacio Muestral):}
\[ \binom{8}{2} = \frac{8!}{2!(8-2)!} = \frac{8 \times 7 \times 6!}{2 \times 1 \times 6!} = \frac{56}{2} = 28 \]

\textbf{Procedimiento de sustitución y cálculo de cada término:}
Se aplica la fórmula de combinatoria $\frac{n!}{k!(n-k)!}$ para cada término del numerador:

\begin{itemize}
    \item \textbf{Cálculo de $f(0,0)$:}
    \[ f(0,0) = \frac{\left[ \frac{3!}{0!3!} \right] \times \left[ \frac{2!}{0!2!} \right] \times \left[ \frac{3!}{2!1!} \right]}{28} = \frac{(1) \times (1) \times (3)}{28} = \frac{3}{28} \]
    
    \item \textbf{Cálculo de $f(1,0)$:}
    \[ f(1,0) = \frac{\left[ \frac{3!}{1!2!} \right] \times \left[ \frac{2!}{0!2!} \right] \times \left[ \frac{3!}{1!2!} \right]}{28} = \frac{(3) \times (1) \times (3)}{28} = \frac{9}{28} \]
    
    \item \textbf{Cálculo de $f(0,1)$:}
    \[ f(0,1) = \frac{\left[ \frac{3!}{0!3!} \right] \times \left[ \frac{2!}{1!1!} \right] \times \left[ \frac{3!}{1!2!} \right]}{28} = \frac{(1) \times (2) \times (3)}{28} = \frac{6}{28} \]
    
    \item \textbf{Cálculo de $f(1,1)$:}
    \[ f(1,1) = \frac{\left[ \frac{3!}{1!2!} \right] \times \left[ \frac{2!}{1!1!} \right] \times \left[ \frac{3!}{0!3!} \right]}{28} = \frac{(3) \times (2) \times (1)}{28} = \frac{6}{28} \]
\end{itemize}

\textbf{Suma Final de Probabilidades:}
\[ P(X \le 1, Y \le 1) = \frac{3}{28} + \frac{9}{28} + \frac{6}{28} + \frac{6}{28} = \frac{24}{28} = \mathbf{0.8571} \]

\section{Aplicación 2: Caso Continuo (Uso de Recursos)}
\textbf{Enunciado:} Función de densidad conjunta $f(x,y) = x + y$ para $0 \le x, y \le 1$. Hallar la probabilidad de que el uso de ambos recursos no supere el 50\%, es decir, $P(X \le 0.5, Y \le 0.5)$.

\textbf{Planteamiento de la Integral Doble:}
\[ P = \int_{0}^{0.5} \int_{0}^{0.5} (x + y) \, dy \, dx \]

\textbf{Desarrollo de la Integración Paso a Paso:}
1. Integral interna respecto a $y$:
\[ \int_{0}^{0.5} (x + y) \, dy = \left[ xy + \frac{y^2}{2} \right]_{y=0}^{y=0.5} = (0.5x + 0.125) \]
2. Integral externa respecto a $x$:
\[ \int_{0}^{0.5} (0.5x + 0.125) \, dx = \left[ \frac{0.5x^2}{2} + 0.125x \right]_{x=0}^{x=0.5} \]
Sustitución de los límites:
\[ P = \left( \frac{0.5(0.5)^2}{2} + 0.125(0.5) \right) = 0.0625 + 0.0625 = \mathbf{0.125} \]

\end{document}
